% ============================================================================
% Fichier  : partie4/P04_C23_audit_production_continuite.tex
% Titre    : Auditer la Production et la Continuité
% Auteur   : MINKA MI NGUIDJOÏ Thierry Emmanuel
% Version  : 1.0 -- novembre 2026
% ============================================================================
% Base : ACSI C17, sections production, sauvegardes et continuité (PX.1 à
% PX.5 ; SANAGA DMIA 8 h, RTO 4 h, RPO 15 min, plan 2022 jamais testé ;
% panne du 9 mars 7 h 30 au 12 mars 8 h, 72,5 h ; dix contrôles 2/7/1).
% Données du journal du premier trimestre 2026 (base des labs ACSI-LAB-09
% et ASI-LAB-13) : 90 jours, 88 lignes, 84 succès, absences les 14 février
% et 8 mars, un échec sans ticket le 27 février ; restaurations 2025 :
% deux tests sur quatre, 3 h 40 et 5 h 10 ; site de secours de Bonabéri ;
% plan approuvé le 12 mai 2022 ; ancien système SIBANK 4.
% Sources : guide du CHAI (fiches 3.3 et 3.5) ; COBIT 2019 (DSS01, DSS02,
% DSS04, BAI09) ; ISO 22301:2019 ; COBAC R-2016/04 (risque opérationnel).
% Corrigés : corriges/Corr_C23.tex.
% ============================================================================

\chapter{Auditer la Production et la Continuité}
\label{chap:audit-continuite}
\chaptermark{Production et continuité}

\epigraph{%
    « Une sauvegarde qu'on n'a jamais restaurée\\
    est une hypothèse, pas une sauvegarde. »%
}{Adage des exploitants}

\niveauChapitre{Intermédiaire}{%
    Maîtrise du Chapitre~\ref{chap:audit-technique} (infrastructure) et du
    Chapitre~\ref{chap:preuve} (tests sur population)%
}

\begin{boxobjectifs}
\begin{itemize}
    \item Décomposer l'audit de la production~: exploitation, sauvegardes et
          restaurations, support, parc, maintenance.
    \item Auditer des sauvegardes par un rapprochement jour par jour, puis
          par leurs restaurations.
    \item Calculer et interpréter les écarts entre RTO, RPO et durées
          constatées.
    \item Auditer un plan de continuité selon l'ISO~22301~: analyse
          d'impact, objectifs, stratégie, plan, exercices, revue.
    \item Assister à un test de restauration et en tirer une preuve.
\end{itemize}
\end{boxobjectifs}

\bigskip

% ============================================================================
\section*{Question directrice : la banque pouvait-elle reprendre en quatre heures~?}
% ============================================================================

Le plan de continuité de la BANQUE SANAGA, approuvé le 12~mai~2022, fixe
pour le système central une durée maximale d'interruption admissible de
8~heures, un objectif de reprise (RTO) de 4~heures et une perte de données
maximale (RPO) de 15~minutes. Le 9~mars~2026, les agences ouvrent sans
système à 7~h~30~; le service revient le 12~mars à 8~h. Personne n'a ouvert
le plan.

La COBAC demande à Mme~EKAMBI si le dispositif de continuité était
capable, avant la panne, de tenir ses objectifs.

\medskip
\textbf{Plan du chapitre.}
\begin{enumerate}
    \item Décomposer l'audit de la production (section~\ref{sec:c23-production}).
    \item Auditer les sauvegardes (section~\ref{sec:c23-sauvegardes}).
    \item Les objectifs de reprise, en chiffres (section~\ref{sec:c23-objectifs}).
    \item Auditer la continuité selon l'ISO~22301 (section~\ref{sec:c23-continuite}).
    \item Le dispositif de la BANQUE SANAGA (section~\ref{sec:c23-sanaga}).
\end{enumerate}
Les prestataires et le nuage, dont dépend souvent la continuité, font
l'objet du Chapitre~\ref{chap:audit-tiers}.

\separateur

% ============================================================================
\section{Décomposer l'Audit de la Production}
\label{sec:c23-production}
% ============================================================================

La production, ou exploitation, fait fonctionner chaque jour les systèmes
que les projets ont livrés. Le guide du CHAI en propose deux fiches, l'une
pour la production, l'autre pour le support et le parc~\cite{chai2014guide}.

\begin{tableaucours}{cL{3.0cm}L{5.6cm}L{2.2cm}}{%
    \tableauHeader{Code} &
    \tableauHeader{Sous-domaine} &
    \tableauHeader{Contrôles (d'après le guide du CHAI)} &
    \tableauHeader{COBIT}
}
PX.1 & Exploitation & Traitements planifiés et suivis~; alertes remontées
en temps réel~; livraisons en production contrôlées par la production~;
guides d'exploitation à jour & DSS01 \\
PX.2 & Sauvegardes et restaurations & Sauvegardes conformes aux besoins~;
comptes rendus contrôlés~; supports externalisés~; restaurations testées
selon un planning & DSS04 \\
PX.3 & Support & Centre de services~; tickets qualifiés et escaladés~;
délais de clôture suivis~; satisfaction mesurée & DSS02 \\
PX.4 & Parc & Inventaire du matériel et des logiciels~; licences~;
installation à partir d'une configuration standard & BAI09 \\
PX.5 & Maintenance & Contrats de maintenance des matériels critiques et des
logiciels de base~; obsolescence suivie & DSS01, BAI09 \\
\end{tableaucours}
\vspace{-0.8em}
\begin{center}{\small\itshape Tableau~23.1. Décomposition de l'audit de la
production}\end{center}

Chez LiZbiZ, les domaines~B (sauvegardes) et~E (exploitation) du
Chapitre~\ref{chap:conclure} relèvent de cette décomposition~: le
domaine~E, en non-conformité opaque ($U = 0{,}30$, puis $0{,}40$ au second
cycle), appelle d'abord la recherche des preuves manquantes, notamment les
comptes rendus des traitements de nuit.

\begin{boxprocedure}
\textbf{Auditer les traitements de nuit (PX.1)}
\begin{enumerate}
    \item Obtenir le plan d'ordonnancement (liste des traitements, ordre,
          heures, dépendances).
    \item Obtenir les comptes rendus d'exécution de la période.
    \item Rapprocher, nuit par nuit, le planifié et l'exécuté~; repérer les
          traitements absents, en échec ou relancés.
    \item Pour chaque échec, retrouver le ticket et l'action~; vérifier
          qu'une relance manuelle a été validée.
    \item Vérifier qu'une personne contrôle les comptes rendus chaque matin
          (visa, ticket, alerte).
\end{enumerate}
\end{boxprocedure}

% ============================================================================
\section{Auditer les Sauvegardes}
\label{sec:c23-sauvegardes}
% ============================================================================

\subsection{Le journal ne suffit pas}

Un journal de sauvegarde dit qu'une copie a été écrite. Il ne dit pas
qu'elle est complète, lisible, ni qu'on saura la remonter dans le délai
voulu. Le guide du CHAI demande donc à l'auditeur de croiser les
sauvegardes avec les besoins, de vérifier le suivi des comptes rendus et
des supports, puis de se concentrer sur la restauration~: procédures
écrites, planning de tests, comptes rendus et actions qui en découlent.

\begin{tableaucours}{L{3.6cm}L{4.4cm}L{4.2cm}}{%
    \tableauHeader{Question} &
    \tableauHeader{Preuve} &
    \tableauHeader{Test}
}
Toutes les sauvegardes prévues ont-elles eu lieu~? & Journal de
sauvegarde, plan de sauvegarde & Rapprocher jour par jour le planifié et
l'exécuté \\
Les échecs sont-ils traités~? & Tickets, comptes rendus & Pour chaque échec,
retrouver l'action et sa date \\
Les supports sont-ils externalisés~? & Bordereaux de sortie & Rapprocher
les dates d'externalisation \\
Sait-on restaurer, et en combien de temps~? & Comptes rendus de
restauration & Comparer la durée constatée au RTO \\
\end{tableaucours}
\vspace{-0.8em}
\begin{center}{\small\itshape Tableau~23.2. Auditer les sauvegardes}\end{center}

\subsection{Le rapprochement jour par jour}

\begin{boxoutils}{Rapprocher le planifié et l'exécuté}
\begin{lstlisting}[language=SQL]
-- calendrier(jour) : tous les jours de la période
-- journal(jour, statut, ticket) : une ligne par sauvegarde exécutée
SELECT c.jour,
       COALESCE(j.statut, 'ABSENTE') AS statut,
       j.ticket
FROM   calendrier c
LEFT JOIN journal j ON j.jour = c.jour
WHERE  j.statut IS NULL OR j.statut <> 'SUCCES'
ORDER  BY c.jour;
\end{lstlisting}
La jointure externe fait apparaître les jours \emph{sans} ligne de
journal~: une sauvegarde absente n'a, par définition, laissé aucune trace
d'échec.
\end{boxoutils}

\begin{boxscenario}{Opération SANAGA. Le premier trimestre 2026}
\textbf{Données.} 90~jours du 1\textsuperscript{er}~janvier au
31~mars~2026~; 88~lignes de journal~; 84~sauvegardes en succès.

\medskip
\textbf{Constatations.} Deux jours sans ligne (14~février et 8~mars)~;
quatre échecs, dont un, le 27~février, sans ticket. Taux de réussite~:
$84/90 \approx 93~\%$, pour un seuil de 98~\% fixé par la procédure.

\medskip
\textbf{Analyse} (confiance élevée). Le contrôle « sauvegardes quotidiennes
réussies » est \NON{}. La sauvegarde du 8~mars, absente, est celle de la
veille de la panne~: le RPO de 15~minutes n'aurait de toute façon pas pu
être tenu par la seule sauvegarde quotidienne.
\end{boxscenario}

\begin{boxalert}
\textbf{Violation typique.} Conclure que les sauvegardes sont efficaces
parce que le journal affiche « succès » chaque nuit.

\smallskip
\textbf{Règle~:} une sauvegarde est efficace quand une restauration
récente, documentée, a réussi dans le délai visé. Sans test de
restauration, le contrôle « on sait restaurer » est \NEK{}, quel que soit
le journal.
\end{boxalert}

% ============================================================================
\section{Les Objectifs de Reprise, en Chiffres}
\label{sec:c23-objectifs}
% ============================================================================

\begin{boxdef}
\termecle{DMIA}\index{durée maximale d'interruption admissible}~: durée au
delà de laquelle l'interruption d'un service met en danger l'organisation.

\medskip
\termecle{RTO}\index{RTO} (\textit{recovery time objective})~: durée visée
pour rétablir le service après une interruption~; il doit être inférieur à
la DMIA.

\medskip
\termecle{RPO}\index{RPO} (\textit{recovery point objective})~: perte de
données maximale admise, exprimée en durée~: les données des dernières
minutes ou heures qui peuvent être perdues.
\end{boxdef}

\subsection{Mesurer les écarts}

\begin{tableaucours}{L{4.6cm}cccc}{%
    \tableauHeader{Événement} &
    \tableauHeader{Objectif} &
    \tableauHeader{Constaté} &
    \tableauHeader{Écart} &
    \tableauHeader{Cote}
}
Restauration de test, 2025 (1) & RTO 4~h & 3~h~40 & $-20$~min & \OUI{} \\
Restauration de test, 2025 (2) & RTO 4~h & 5~h~10 & $+70$~min & \NON{} \\
Panne de mars~2026 & RTO 4~h & 72~h~30 & $+68$~h~30 & \NON{} \\
Panne de mars~2026 & DMIA 8~h & 72~h~30 & $\times 9$ & \NON{} \\
\end{tableaucours}
\vspace{-0.8em}
\begin{center}{\small\itshape Tableau~23.3. Objectifs et durées constatées,
BANQUE SANAGA}\end{center}

\begin{boxinfo}
\textbf{Le RPO ne se juge pas avec une sauvegarde quotidienne.} Un RPO de
15~minutes suppose une réplication continue (ou des journaux de
transactions copiés toutes les quinze minutes au plus) vers le site de
secours. L'auditeur vérifie le dispositif de réplication, son retard
constaté (en minutes) et ses alertes~; une restauration de la sauvegarde de
la nuit ne prouve rien sur le RPO.
\end{boxinfo}

% ============================================================================
\section{Auditer la Continuité selon l'ISO 22301}
\label{sec:c23-continuite}
% ============================================================================

\subsection{La démarche}

L'ISO~22301 formule les exigences d'un système de management de la
continuité d'activité, sur la même structure que l'ISO/IEC~27001 (clauses
4 à 10)~\cite{iso22301_2019}. Son cœur opérationnel (clause~8) enchaîne
l'analyse d'impact, l'appréciation des risques, la stratégie de
continuité, les plans, puis les exercices. Le cours de gouvernance
(MSN-2S3-MSI03) en présente la conception~; l'auditeur suit la même chaîne,
et cherche à chaque maillon la pièce qui le prouve.

\begin{boxprocedure}
\textbf{Auditer un plan de continuité}
\begin{enumerate}
    \item \textbf{Analyse d'impact}~: datée, revue après chaque changement
          majeur, approuvée par les métiers.
    \item \textbf{Objectifs}~: DMIA, RTO et RPO fixés pour chaque service
          critique, et cohérents (RTO inférieur à la DMIA).
    \item \textbf{Stratégie et moyens}~: le site de secours héberge la
          version en production, et les données y sont répliquées au rythme
          du RPO.
    \item \textbf{Plan}~: il décrit le système \emph{actuel}, les
          personnes à joindre, les procédures de bascule et de retour.
    \item \textbf{Exercices}~: au moins une fois par an pour chaque service
          critique, avec un compte rendu, la durée constatée et les actions.
    \item \textbf{Écart et revue}~: la durée constatée, en exercice ou lors
          d'un incident réel, est comparée au RTO~; une revue après incident
          est formalisée.
\end{enumerate}
\end{boxprocedure}

\subsection{Quatre types d'exercices}

\begin{tableaucours}{L{3.2cm}L{5.2cm}L{3.8cm}}{%
    \tableauHeader{Exercice} &
    \tableauHeader{Principe} &
    \tableauHeader{Ce qu'il prouve}
}
Revue sur table & Les responsables déroulent le plan sur un scénario,
autour d'une table & Que le plan est compris et cohérent \\
Test technique partiel & On restaure un système sur le site de secours &
Que la restauration est possible, et sa durée \\
Bascule réelle planifiée & On fait tourner un service depuis le site de
secours & Que la reprise tient le RTO \\
Exercice inopiné & Déclenchement sans préavis & Que l'organisation réagit
sans préparation \\
\end{tableaucours}
\vspace{-0.8em}
\begin{center}{\small\itshape Tableau~23.4. Les exercices de continuité}\end{center}

Une revue sur table ne prouve pas que le RTO est tenable~; seule une
bascule réelle le prouve. L'auditeur cote donc le contrôle « RTO tenu en
exercice » à partir de la nature de l'exercice.

% ============================================================================
\section{Le Dispositif de la BANQUE SANAGA}
\label{sec:c23-sanaga}
% ============================================================================

\begin{boxlecture}{Continuité, système central}
Mme~EKAMBI cote dix contrôles de continuité du système central, à partir
du plan de 2022, des journaux de sauvegarde du premier trimestre~2026, des
comptes rendus de restauration de 2025 et de la chronologie de la panne.

\medskip
\begin{tableaucours}{L{7.0cm}c}{%
    \tableauHeader{Contrôle} &
    \tableauHeader{Cote}
}
Analyse d'impact revue depuis 2022 & \NON{} \\
RTO et RPO fixés et approuvés & \OUI{} \\
Plan à jour du système central actuel & \NON{} \\
Site de secours à la version de production & \NON{} \\
Plan testé dans l'année & \NON{} \\
RTO tenu lors du dernier incident & \NON{} \\
Sauvegardes quotidiennes réussies (seuil~: 98~\%) & \NON{} \\
Supports de sauvegarde externalisés & \OUI{} \\
Restauration testée chaque trimestre & \NON{} \\
Revue après la reprise formalisée & \NEK{} \\
\end{tableaucours}

\medskip
\textbf{Constatation.} $\mu = 0{,}20$~; $U = 0{,}10$~; $H \approx 1{,}16$~bit.
Le service a été indisponible 72,5~heures pour un RTO de 4~heures. Sur les
90~jours du premier trimestre, 84~sauvegardes ont réussi (93~\%). En 2025,
deux des quatre tests de restauration prévus ont eu lieu, dont un a
dépassé le RTO.

\medskip
\textbf{Analyse} (confiance élevée). Non-conformité claire~: la banque ne
pouvait pas tenir ses objectifs. Les objectifs existaient, les moyens ne
suivaient plus le système~: le plan et le site de secours de Bonabéri
décrivaient l'ancien système central (SIBANK~4). Le seul \NEK{} porte sur
la revue après reprise~: le rapport à la COBAC existe, mais aucune revue
interne formalisée n'a été retrouvée.

\medskip
\hfill\lecture{0{,}20}{0{,}10}{1{,}16}
\end{boxlecture}

\begin{boxmetier}{Auditeur SI}
\textbf{Assister à un test de restauration}
\begin{enumerate}
    \item Choisir soi-même la sauvegarde à restaurer, et la date, dans la
          période auditée.
    \item Noter l'heure de début, les incidents, l'heure à laquelle le
          service est utilisable.
    \item Faire vérifier par un utilisateur métier que les données
          restaurées sont complètes (quelques opérations connues).
    \item Comparer la durée au RTO, et la date des dernières données
          restaurées au RPO.
    \item Conserver le compte rendu, signé, au dossier.
\end{enumerate}
\end{boxmetier}

\begin{boxmetier}{Gestionnaire de risques}
\textbf{Relier continuité et changements.} La leçon de la BANQUE SANAGA
est qu'un changement majeur (la migration) a rendu le plan obsolète sans
que personne ne le voie. Le gestionnaire de risques inscrit au registre une
règle~: tout changement majeur déclenche la revue de l'analyse d'impact, du
plan et du site de secours, avant la mise en production.
\end{boxmetier}

\begin{boxinfo}
\textbf{Ce que dit la recherche.} L'axiome d'individuation temporelle
d'A0~\cite{minka2026auditologie} rappelle que l'objet audité change~: un
plan de continuité exact en 2022 peut être faux en 2026 sans qu'aucune de
ses lignes n'ait changé. L'auditeur date donc chaque preuve, et vérifie
qu'elle décrit le système de la période auditée.
\end{boxinfo}

\separateur

\begin{boxresume}
\textbf{Ce qu'il faut retenir de ce chapitre~:}
\begin{itemize}
    \item \textbf{Cinq sous-domaines de production}~: exploitation,
          sauvegardes et restaurations, support, parc, maintenance.
    \item \textbf{Les sauvegardes} se rapprochent jour par jour (une
          sauvegarde absente ne laisse pas d'échec), puis s'auditent par
          leur restauration.
    \item \textbf{RTO, RPO, DMIA}~: on mesure l'écart entre objectif et
          constaté~; le RPO ne se juge pas avec une sauvegarde quotidienne.
    \item \textbf{La continuité est une chaîne} (ISO~22301)~: analyse
          d'impact, objectifs, stratégie et moyens, plan, exercices, revue.
    \item \textbf{Un plan qui décrit l'ancien système} ne protège pas le
          nouveau~: tout changement majeur déclenche sa revue.
\end{itemize}
\end{boxresume}

\begin{boxinfo}
\textbf{Points de contrôle de l'élève.} Je sais~: décomposer la production
en cinq sous-domaines~; écrire le rapprochement jour par jour~; calculer
les écarts à RTO et DMIA~; dire ce qu'il faut pour juger un RPO~; auditer un
plan de continuité en six maillons~; distinguer les quatre types
d'exercices.
\end{boxinfo}

\bigskip

\begin{boxexo}[title={\ding{45}\quad Exercice~23.1. Rattacher les contrôles \niveaubase}]
Rattachez chaque contrôle à un sous-domaine (PX.1 à PX.5).
\begin{enumerate}[(a)]
    \item Les licences installées sont rapprochées des licences acquises.
    \item Les traitements de nuit ont un compte rendu contrôlé chaque matin.
    \item Le serveur de base de données est couvert par un contrat de
          maintenance avec intervention en quatre heures.
    \item Les tickets de priorité 1 sont escaladés en moins de 30~minutes.
    \item Un test de restauration est mené chaque trimestre.
\end{enumerate}
\end{boxexo}

\begin{boxexo}[title={\ding{45}\quad Exercice~23.2. RTO et RPO \niveaumoyen}]
Lors d'un test de restauration du système central de la BANQUE SANAGA, la
restauration commence à 9~h~00 et le service est utilisable à 14~h~10. La
dernière opération présente dans les données restaurées date de la veille
à 23~h~50~; la sauvegarde testée avait été lancée à 23~h~30.
\begin{enumerate}[(a)]
    \item Le RTO de 4~heures est-il tenu~? Calculez l'écart.
    \item Pourquoi ce test ne permet-il pas de juger le RPO de
          15~minutes~?
    \item Quel dispositif le RPO de 15~minutes suppose-t-il, et comment
          l'auditer~?
    \item Cotez le contrôle « RTO tenu en test » et rédigez la
          constatation.
\end{enumerate}
\end{boxexo}

\begin{boxexo}[title={\ding{45}\quad Exercice~23.3. Le rapprochement \niveaumoyen}]
Exécutez à la main la requête de rapprochement sur la semaine du 6 au
12~mars~2026, sachant que le journal contient des succès les 6, 7, 9, 10,
11 et 12, et rien le 8. Que rend la requête~? Pourquoi la jointure externe
est-elle indispensable~?
\end{boxexo}

\begin{boxexo}[title={\ding{45}\quad Exercice~23.4. Le taux de réussite \niveaubase}]
Recalculez le taux de réussite du premier trimestre de la BANQUE SANAGA
(a) en comptant les deux jours absents comme des échecs, (b) en ne
comptant que les lignes du journal. Lequel est juste, et pourquoi~?
\end{boxexo}

\begin{boxexo}[title={\ding{45}\quad Exercice~23.5. Quel exercice~? \niveaumoyen}]
Le plan de continuité de LiZbiZ prévoit « un exercice annuel ». Le compte
rendu de 2025 décrit une réunion de deux heures où les responsables ont lu
le plan. Le contrôle « RTO tenu en exercice » est-il \OUI{}, \NON{} ou
\NEK{}~? Justifiez.
\end{boxexo}

\begin{boxexo}[title={\ding{45}\quad Exercice~23.6. Les traitements de nuit \niveaumoyen}]
Le plan d'ordonnancement de LiZbiZ prévoit 12~traitements chaque nuit. Sur
30~nuits, les comptes rendus montrent 348~exécutions réussies, 6~échecs
relancés le matin sans validation, et 6~absences. Calculez les taux,
rédigez la constatation et cotez PX.1.
\end{boxexo}

\begin{boxexo}[title={\ding{45}\quad Exercice~23.7. La revue après incident \niveaumoyen}]
Rédigez la fiche d'assertion du contrôle « une revue après reprise est
formalisée », et dites ce qui permettrait à Mme~EKAMBI de lever le \NEK{}.
\end{boxexo}

\begin{boxexo}[title={\ding{45}\quad Exercice~23.8. Les dix contrôles \niveaumoyen}]
Recalculez $H$ pour la lecture de la continuité de la BANQUE SANAGA (2/7/1)
et décomposez-la en ignorance et en désaccord. Que montre la
décomposition~?
\end{boxexo}

\begin{boxexo}[title={\ding{45}\quad Exercice~23.9. Le rapport à la COBAC \niveauexpert}]
Rédigez l'observation en sept éléments relative au plan de continuité
obsolète, et la réponse de Mme~EKAMBI à la question de la COBAC (« le
dispositif était-il capable de tenir ses objectifs~? »), en une page.
\end{boxexo}

\begin{boxexo}[title={\ding{45}\quad Exercice~23.10. Un programme d'exercices \niveauexpert}]
Proposez à la BANQUE SANAGA un programme d'exercices de continuité sur deux
ans, pour le système central, la messagerie et la monétique, avec le type
d'exercice, la fréquence, ce qu'il prouve, et la manière dont l'audit
interne en utilisera les comptes rendus.
\end{boxexo}

\bigskip

\begin{boxinfo}
\textbf{Pour aller plus loin.}
\begin{itemize}
    \item ISO~22301:2019~\cite{iso22301_2019}~: le système de management de
          la continuité.
    \item Le guide du CHAI~\cite{chai2014guide}~: les fiches
          « production » et « support et parc ».
    \item ISACA, COBIT~2019~\cite{isaca2018cobitobjectifs}~: DSS01, DSS02
          et DSS04.
\end{itemize}
\end{boxinfo}

\begin{boxinfo}
\textbf{Transition.} La continuité de la BANQUE SANAGA dépend aussi de son
prestataire, DIGICORE AFRIQUE. Le Chapitre~\ref{chap:audit-tiers} audite
les tiers et le nuage~: contrats, droit d'audit, rapports d'assurance,
responsabilité partagée, localisation des données, réversibilité.
\end{boxinfo}
